\section{Empirical Framework}
\label{sec:framework}

\subsection{Separating pricing and sourcing margins}
\label{subsec:pricing_sourcing}

Let a purchase price reflect both the price charged by the winning firm and the process that determines which firm wins. A court order can raise costs through a same-firm markup, through reduced scale, through lower participation, or through a different supplier match. The empirical framework separates these channels with three objects: urgent-versus-ordinary estimates, administrative-versus-litigated urgent estimates, and within firm-buyer-item comparisons.

The within-firm test is the pricing test. It estimates whether the same supplier charges a different price to the same buyer for the same item when the purchase is administrative rather than litigated. Because this comparison conditions on the winning firm, it cannot measure supplier reallocation. That is the point: if the overall gap remains while the within-firm gap disappears, the cost margin is not a broad same-firm markup.

\subsection{Urgent-vs-ordinary procurement}
\label{subsec:uo}

The first specification compares urgent and ordinary purchases:
\begin{equation}
y_{igbt}=\beta U_{igbt}+\alpha_i+\gamma_t+\delta_b+\varepsilon_{igbt},
\label{eq:urgent_ordinary}
\end{equation}
where $y_{igbt}$ is a procurement outcome for item $i$, supplier or purchase observation $g$, buyer $b$, and year $t$; $U_{igbt}$ indicates urgent procurement; and the fixed effects absorb item, year, and buyer differences. This estimate is descriptive of the urgent cost margin, not the core sanction estimate.

\subsection{Administrative-vs-litigated urgent procurement}
\label{subsec:utg}

Within urgent procurement, we estimate
\begin{equation}
y_{igbt}=\theta A_{igbt}+\alpha_i+\gamma_t+\delta_b+\varepsilon_{igbt},
\label{eq:utg}
\end{equation}
where $A_{igbt}$ equals one for administrative urgent purchases and zero for litigated urgent purchases. Negative estimates in price regressions mean administrative purchases are cheaper, or equivalently that litigated purchases are more expensive. Because administrative purchases are selected and larger, $\theta$ is not interpreted alone. The paper reports selection bounds and decomposes the scale channel.

\subsection{Selection bounds}
\label{subsec:bounds}

Administrative requests are screened by the health agency. If screened cases are more feasible or less costly, naive comparisons overstate the sanction-related gap. We therefore use Manski-Lee trimming bounds. In strata where administrative observations exceed litigated observations, we trim the administrative price distribution by the differential selection proportion. Trimming the high tail and low tail gives bounds on the litigated-over-administrative gap. The preferred bounds are reported as percentage price gaps to avoid sign ambiguity.

The bounds require monotone selection into the administrative channel within trimming strata. This is weaker than random assignment but still substantive: it assumes that the screened administrative cases form an ordered subset relative to the potential litigated cases within the stratum. Online Appendix C gives the formal trimming logic and reports the trimming proportions.
